% Handschrift: Teil 7 - Komplexe Zahlen, Seiten 1--8.
\section{Von den reellen Zahlen zur komplexen Zahlenebene}
\label{sec:complex_numbers_cartesian_form}

Komplexe Zahlen vereinfachen viele Rechnungen und treten in der Beschreibung
physikalischer Zusammenhänge auf. Zunächst motiviert uns eine Gleichung,
deren Koeffizienten reell sind, die aber keine reelle Lösung besitzt:
\[
  z^2-6z+25\overset{!}{=}0
  \quad\Longleftrightarrow\quad (z-3)^2=-16.
\]
Auch die $pq$-Formel führt auf $3\pm\sqrt{9-25}$ und damit auf eine negative
Zahl unter der Wurzel. Können wir den Zahlbereich so erweitern, dass diese
und schließlich alle nichtkonstanten Polynomgleichungen lösbar werden?

\par\noindent\begin{minipage}{\linewidth}
\begin{DefinitionBox}{Imaginäre Einheit und komplexe Zahlen}{complex_numbers_cartesian}
  Die \emph{imaginäre Einheit} \EnglishTerm{imaginary unit} $i$ erfüllt
  $i^2=-1$. Da keine reelle Zahl diese Eigenschaft hat, gilt $i\notin\R$.

  Eine \emph{komplexe Zahl} \EnglishTerm{complex number} besitzt die
  \emph{kartesische Darstellung} \EnglishTerm{Cartesian form}
  \[
    z=x+iy,\qquad x,y\in\R,\qquad
    \C=\{x+iy\mid x,y\in\R\}.
  \]
  Dabei heißen $\operatorname{Re}(z)=x$ der \emph{Realteil}
  \EnglishTerm{real part} und $\operatorname{Im}(z)=y$ der
  \emph{Imaginärteil} \EnglishTerm{imaginary part}.
\end{DefinitionBox}
\vspace{-0.5\baselineskip}
\begin{ExampleBox}{Realteil und Imaginärteil ablesen}{complex_numbers_components}
  Bestimmen Sie die beiden reellen Komponenten der folgenden Zahlen.
  \begin{center}
    \renewcommand{\arraystretch}{1.3}
    \begin{tabular}{@{}c@{\qquad}c@{\qquad}c@{}}
      $z$ & $\operatorname{Re}(z)$ & $\operatorname{Im}(z)$\\\midrule
      $2+3i$ & $\Cloze{2cm}{2}$ & $\Cloze{2cm}{3}$\\
      $i$ & $\Cloze{2cm}{0}$ & $\Cloze{2cm}{1}$\\
      $3$ & $\Cloze{2cm}{3}$ & $\Cloze{2cm}{0}$
    \end{tabular}
  \end{center}
\end{ExampleBox}
\end{minipage}\par


\begin{ExampleBox}{Die Ausgangsgleichung lösen}{complex_numbers_initial_equation}
  Bestimmen Sie mithilfe von $i^2=-1$ die beiden Lösungen der Ausgangsgleichung.
  \[
    (z-3)^2=(4i)^2
      \quad\Longleftrightarrow\quad z=\Cloze{3cm}{3\pm4i}.
  \]
  Die Probe bestätigt $(\pm4i)^2=16i^2=-16$.
\end{ExampleBox}

\clearpage
\begin{RemarkBox}[remark title=Realteil und Imaginärteil sind reell]
  Realteil und Imaginärteil sind beide \emph{reell}; der Imaginärteil von
  $2+3i$ ist $3$, nicht $3i$.
  Zahlen mit Imaginärteil $0$ sind reell, Zahlen mit Realteil $0$ heißen
  \emph{rein imaginär} \EnglishTerm{purely imaginary}. Damit gilt
  \[
    z\in\R\ \Longleftrightarrow\ \operatorname{Im}(z)=0,
    \qquad \N\subsetneq\Z\subsetneq\Q\subsetneq\R\subsetneq\C.
  \]
\end{RemarkBox}

\subsection{Zahlenebene und Betrag}
Die \emph{Gaußsche Zahlenebene} \EnglishTerm{complex plane} ordnet
$z=x+iy$ dem Punkt $(x,y)$ zu. Die waagerechte \emph{reelle Achse}
\EnglishTerm{real axis} trägt den Realteil, die senkrechte
\emph{imaginäre Achse} \EnglishTerm{imaginary axis} den Imaginärteil.
Jede reelle Zahl liegt auf der reellen Achse, jede rein imaginäre Zahl
auf der imaginären Achse.
\begin{center}
  \CachedTikzFigure{content/resources/chapters/complex_numbers/tikz/complex_plane_cartesian.tex}
\end{center}
Der Pfeil vom Ursprung zum Punkt veranschaulicht den Abstand vom Ursprung.
Aus dem rechtwinkligen Dreieck folgt mit dem Satz des Pythagoras die Betragsformel.

\par\noindent\begin{minipage}{\linewidth}
\begin{DefinitionBox}{Betrag einer komplexen Zahl}{complex_numbers_modulus}
  Der \emph{Betrag} \EnglishTerm{modulus} einer komplexen Zahl $z=x+iy$
  ist ihr Abstand vom Ursprung:
  \[
    |z|=\sqrt{x^2+y^2}
       =\sqrt{\bigl(\operatorname{Re}(z)\bigr)^2+
                    \bigl(\operatorname{Im}(z)\bigr)^2}\in\R_{\ge0}.
  \]
\end{DefinitionBox}
\vspace{-0.5\baselineskip}
\begin{ExampleBox}{Eine Zahl als Punkt und als Pfeil}{complex_numbers_plane_example}
  Bestimmen Sie den Betrag von $z=3+2i$ anhand der eingezeichneten Komponenten.
  \par\smallskip
  \begin{minipage}[c]{0.50\linewidth}
    \CachedTikzFigure{content/resources/chapters/complex_numbers/tikz/complex_plane_modulus.tex}
  \end{minipage}\hfill
  \begin{minipage}[c]{0.45\linewidth}
    Die Katheten haben die Längen $3$ und $2$. Der Abstand beträgt daher
    \[
      |3+2i|=\Cloze{3cm}{\sqrt{3^2+2^2}=\sqrt{13}}.
    \]
  \end{minipage}
\end{ExampleBox}
\end{minipage}\par

\clearpage
\subsection{Komplexe Konjugation}
Zur Konjugation betrachten wir die Spiegelung eines Punktes an der reellen Achse.

\par\noindent\begin{minipage}{\linewidth}
\begin{DefinitionBox}{Konjugiert komplexe Zahl}{complex_numbers_conjugate}
  Zu $z=x+iy$ heißt $z^*=x-iy$ die \emph{konjugiert komplexe Zahl}
  \EnglishTerm{complex conjugate}. Sie entsteht durch
  \emph{Spiegelung an der reellen Achse}.
  \par\medskip
  \begin{minipage}[c]{0.51\linewidth}
    \CachedTikzFigure{content/resources/chapters/complex_numbers/tikz/complex_conjugation.tex}
  \end{minipage}\hfill
  \begin{minipage}[c]{0.44\linewidth}
    \textbf{Einfach jedes $i$ durch $-i$ ersetzen.}

    Der Realteil bleibt unverändert; der Imaginärteil wechselt sein Vorzeichen:
    \[
      \begin{aligned}
        \operatorname{Re}(z^*)&=\operatorname{Re}(z),\\[4pt]
        \operatorname{Im}(z^*)&=-\operatorname{Im}(z).
      \end{aligned}
    \]
    Diese Regel gilt auch in der später eingeführten Exponentialdarstellung.
  \end{minipage}
\end{DefinitionBox}
\vspace{-0.5\baselineskip}
\begin{ExampleBox}{Konjugation und Betrag}{complex_numbers_conjugation_example}
  Bestimmen Sie zu $z=3+2i$ die konjugierte Zahl und vergleichen Sie die Beträge.
  \[
    z^*=\Cloze{3cm}{3-2i},\qquad
    |z|=\sqrt{3^2+2^2}=\sqrt{13},\qquad
    |z^*|=\Cloze{3cm}{\sqrt{3^2+(-2)^2}=\sqrt{13}}.
  \]
\end{ExampleBox}
\end{minipage}\par

\begin{KeyPointBox}{Eigenschaften der Konjugation}{complex_numbers_conjugation_properties}
  Zweimaliges Spiegeln führt zurück zur Ausgangszahl; Punkte auf der
  Spiegelachse bleiben fest. Daher gilt
  \[
    (z^*)^*=z,\qquad
    z=z^*\ \Longleftrightarrow\ \operatorname{Im}(z)=0
          \ \Longleftrightarrow\ z\in\R.
  \]
  Die Spiegelung ändert den Abstand vom Ursprung nicht:
  \[
    |z^*|=\sqrt{x^2+(-y)^2}=\sqrt{x^2+y^2}=|z|.
  \]
\end{KeyPointBox}
